\documentclass[11pt,a4paper]{article}

\usepackage{amsmath,amssymb,amsthm}
\usepackage[utf8]{inputenc}
\usepackage{mathrsfs}
\usepackage[margin=1in]{geometry}
\usepackage{hyperref}
\usepackage{enumitem}

% Theorem environments
\theoremstyle{plain}
\newtheorem{theorem}{Theorem}[section]
\newtheorem{proposition}[theorem]{Proposition}
\newtheorem{lemma}[theorem]{Lemma}
\newtheorem{corollary}[theorem]{Corollary}

\theoremstyle{definition}
\newtheorem{definition}[theorem]{Definition}
\newtheorem{example}[theorem]{Example}

\theoremstyle{remark}
\newtheorem{remark}[theorem]{Remark}

% Operators and shortcuts
\newcommand{\E}{\mathbb{E}}
\newcommand{\R}{\mathbb{R}}
\newcommand{\Z}{\mathbb{Z}}
\newcommand{\C}{\mathbb{C}}
\newcommand{\N}{\mathbb{N}}
\newcommand{\T}{\mathbb{T}}
\newcommand{\HH}{\mathcal{H}_H}
\newcommand{\Id}{\mathrm{Id}}
\newcommand{\Var}{\mathrm{Var}}
\newcommand{\tr}{\mathrm{tr}}

\title{The Iterated $\Gamma$-Calculus for the Fractional Laplacian, II:\\
Operator Realization and Wiener Chaos}

\author{Ramiro Fontes\\
Quijotic Research\\
\texttt{ramirofontes@gmail.com}}

\date{March 2026}

\begin{document}

\maketitle

\begin{abstract}
We prove that the Bakry--\'Emery curvature theory for the
fractional Laplacian, developed spectrally in Part~I,
admits a structural realization through the operator
derivative framework of~\cite{Fontes2026opderiv}.

The operator derivative $D_{B^H} := \delta_{B^H}^*$,
defined by adjointness on the fBM energy space~$\HH$,
satisfies
$\Gamma(f,f)(x) = \|D_{B^H}\Phi_f(x)\|^2_{\HH}$:
the carr\'e du champ of the stable generator is the
energy-space norm of the operator derivative.
The full hierarchy $\Gamma_k \leftrightarrow R_H^{\circ k}$
is realized by a \emph{diagonal tensor lift}
$\iota_k\colon e_n \mapsto e_n^{\otimes k}$
in $\HH^{\otimes k}$, with Gram matrix
$\langle e_n^{\otimes k}, e_m^{\otimes k}\rangle
= R_H(n,m)^k$.

At $\gamma = 1$ ($H = 1/2$), the energy space has
orthonormal increments $\Delta_j = e_j - e_{j-1}$,
and the \emph{increment tensors}
$\delta_j^{(k)} := e_j^{\otimes k} - e_{j-1}^{\otimes k}$
form an orthogonal system in $\HH^{\otimes k}$ with
$\|\delta_j^{(k)}\|^2 = j^k - (j{-}1)^k$.
This gives the Cholesky factorization, the eigenvalue
formula $\kappa_j^{(k)} = \Delta(j^k)/\Delta(j^{k-1})$,
and the curvature stability $\kappa_1^{(k)} = 1$
as consequences of the independent-increment structure
of Brownian motion in the tensor product energy space.
The volume ratio of curvature ellipsoids interpolates
monotonically from $(2N{-}1)!!$ to $N!$ as the chaos
depth increases.

For $\gamma \neq 1$, we prove that $\gamma = 1$ is
the unique global maximizer of the Bakry--\'Emery
curvature, and that the higher-chaos levels regularize
the curvature deficit: $\kappa^{(k)}(\gamma,N) \to 1$
as $k \to \infty$, exponentially in $k$ at rate
$\rho^* = \max_{n \neq m}|\rho_{nm}| < 1$ (the maximum
off-diagonal fBM correlation).

As an analytical consequence, the $k$-th energy
$\mathcal{E}_k(f) = \int \Gamma_k(f,f)\,d\mu$
satisfies the dissipation identity
$\frac{d}{dt}\mathcal{E}_k(P_t f)
= -2\mathcal{E}_{k+1}(P_t f)$
and decays at rate $e^{-2\kappa^{(k+1)}t}$.
By chaos regularization, this rate approaches
$e^{-2t}$---the Gaussian rate---for all
$\gamma \in (0,2)$ at high chaos depth: every
stable semigroup is asymptotically Gaussian
in its high-energy behavior.
Subordination of stable processes acts as a flow on
the Hurst parameter and always decreases curvature.
\end{abstract}

\medskip
\noindent\textbf{MSC 2020:} 60G52, 60G22, 60H07, 35S10, 47D07.\\
\textbf{Keywords:} iterated carr\'e du champ; operator derivative;
energy space; Hadamard power; Wiener chaos;
fractional Brownian motion; Cholesky factorization;
curvature stability; Bakry--\'Emery hierarchy.


%======================================================================
\section{Introduction}\label{sec:intro}
%======================================================================

\subsection{Background}

In Part~I~\cite{Fontes2026b}, we established the correspondence
\begin{equation}\label{eq:part1-duality}
  \Psi_\gamma(\xi,\eta)
  = \tfrac{1}{2}\bigl(|\xi|^\gamma + |\eta|^\gamma
    - |\xi - \eta|^\gamma\bigr)
  = R_{\gamma/2}(|\xi|,|\eta|),
  \qquad \xi\eta \geq 0,
\end{equation}
identifying the Fourier-space carr\'e du champ of the
$\gamma$-stable generator $L_\gamma = -(-\Delta)^{\gamma/2}$
on $\T = \R/(2\pi\Z)$ with the covariance kernel of fractional
Brownian motion (fBM) with Hurst parameter $H = \gamma/2$ on
same-sign frequencies.  This reduced the Bakry--\'Emery curvature
to a generalized eigenvalue problem for fBM covariance matrices,
yielding the Hadamard square identity
$\Gamma_2 \leftrightarrow R_H^{\circ 2}$,
the exact curvature $\kappa(1,N) = 1$ at the Cauchy point
$\gamma = 1$, and the curvature-dimension inequality
$\mathrm{CD}(1-\omega^2/2,\,\infty)$ for the
Cauchy--Fokker--Planck operator with cosine potential.

In a separate work~\cite{Fontes2026opderiv}, we developed the
\emph{operator derivative} $D_X := \delta_X^*$, defined by
Hilbert space adjointness on the energy space $(\mathcal{H}_X,
\delta_X)$ of a process~$X$.  The factorization
$(\Id - \E)\,F = \delta_X(\Pi_X D_X F)$
decomposes $L^2$-fluctuations into differentiation, adaptation,
and integration, yielding a unified Clark--Ocone formula for
all processes with the predictable representation property.

The present paper connects these two frameworks.


\subsection{Main results}

We work on $\T = \R/(2\pi\Z)$ with $\mu = dx/(2\pi)$
and write $\mathcal{T}_N^+$ for the space of
positive-frequency trigonometric polynomials of degree at
most~$N$.  For fBM $B^H$ sampled at $\{1,\ldots,N\}$,
the energy space is $\HH = \R^N$ with inner product
$\langle \mathbf{u}, \mathbf{v} \rangle_{\HH}
= \sum_{n,m} u_n\,v_m\,R_H(n,m)$, so that
$\E[B_n^H\,B_m^H] = \langle e_n, e_m \rangle_{\HH}
= R_H(n,m)$.
The tensor power $\HH^{\otimes k}$ has inner product
$\langle e_{n_1} \otimes \cdots \otimes e_{n_k},\,
e_{m_1} \otimes \cdots \otimes e_{m_k} \rangle
= \prod_{\ell=1}^k R_H(n_\ell, m_\ell)$.

Given $f = \sum_{n=1}^N a_n\,e^{inx} \in \mathcal{T}_N^+$,
define the phase-stripped coefficient vector
$\mathbf{b}(x) = (a_1 e^{ix}, \ldots, a_N e^{iNx})^T$
and the stochastic representative
$\Phi_f(x) := \delta_{B^H}(\mathbf{b}(x))
= \sum_n b_n(x)\,B_n^H$.

\medskip\noindent
\textbf{Theorem~A} (Operator realization of the carr\'e
du champ; \S\ref{sec:realization}).
\emph{For $f \in \mathcal{T}_N^+$ and all $x \in \T$:
\begin{equation}\label{eq:intro-A}
  \Gamma(f,f)(x)
  = \bigl\|D_{B^H}\Phi_f(x)\bigr\|^2_{\HH}.
\end{equation}
The carr\'e du champ of the stable generator on
positive frequencies equals the energy-space norm of the
operator derivative of the stochastic representative.}

\medskip\noindent
\textbf{Theorem~B} (Tensor realization of the hierarchy;
\S\ref{sec:tensor}).
\emph{Define the diagonal tensor lift
$\iota_k\colon \R^N \to \HH^{\otimes k}$ by
$\iota_k(\mathbf{v})
= \sum_n v_n\,e_n^{\otimes k}$.
Then for all $k \geq 1$:
\begin{equation}\label{eq:intro-B}
  \Gamma_k(f,f)(x)
  = \bigl\|\iota_k(\mathbf{b}(x))\bigr\|^2_{\HH^{\otimes k}}.
\end{equation}
The Gram matrix of the diagonal lift is
$\langle e_n^{\otimes k}, e_m^{\otimes k}\rangle_{\HH^{\otimes k}}
= R_H(n,m)^k = R_H^{\circ k}(n,m)$,
and the $k$-th order Bakry--\'Emery curvature is the
worst-case contraction ratio:
$\kappa^{(k)}(\gamma,N) = \lambda_{\min}(R_H^{\circ k},\,
R_H^{\circ(k-1)})$.}

\medskip\noindent
\textbf{Theorem~C} (Increment tensor decomposition;
\S\ref{sec:increment}).
\emph{At $\gamma = 1$ ($H = 1/2$), the energy space
$\mathcal{H}_{1/2}$ has orthonormal increments
$\Delta_j = e_j - e_{j-1}$.
The increment tensors
$\delta_j^{(k)} := e_j^{\otimes k} - e_{j-1}^{\otimes k}$
(with $e_0^{\otimes k} = 0$) satisfy:
\begin{enumerate}[label=\textup{(\roman*)}]
\item Orthogonality:
  $\langle \delta_i^{(k)}, \delta_j^{(k)} \rangle
  = 0$ for $i \neq j$.
\item Norm:
  $\|\delta_j^{(k)}\|^2 = j^k - (j{-}1)^k
  = \Delta(j^k)$.
\item Expansion:
  $e_n^{\otimes k} = \sum_{j=1}^n \delta_j^{(k)}$
  (telescoping).
\end{enumerate}
Consequently, the Cholesky factorization
$\min(n,m)^k = L\,D_k\,L^T$ holds with
$L_{nj} = \mathbf{1}_{n \geq j}$ and
$D_k = \mathrm{diag}(\|\delta_j^{(k)}\|^2)$.}

\medskip\noindent
\textbf{Theorem~D} (Eigenvalue formula;
\S\ref{sec:eigenvalues}).
\emph{At $\gamma = 1$, the $k$-th curvature matrix
$(R_{1/2}^{\circ(k-1)})^{-1}\,R_{1/2}^{\circ k}$
is upper triangular with eigenvalues
\begin{equation}\label{eq:intro-D}
  \kappa_j^{(k)}
  = \frac{\|\delta_j^{(k)}\|^2}{\|\delta_j^{(k-1)}\|^2}
  = \frac{j^k - (j-1)^k}{j^{k-1} - (j-1)^{k-1}},
  \qquad j = 1,\ldots,N.
\end{equation}
In particular:}
\begin{enumerate}[label=\textup{(\roman*)}]
\item \emph{At $k = 2$:
  $\kappa_j^{(2)} = 2j - 1$
  (Part~I, Theorem~C).}
\item \emph{Curvature stability:
  $\kappa_1^{(k)} = 1$ for all $k \geq 2$.}
\item \emph{As $k \to \infty$:
  $\kappa_j^{(k)} \to j$.}
\end{enumerate}

\medskip\noindent
\textbf{Theorem~E} (Volume interpolation;
\S\ref{sec:volume}).
\emph{The determinant of the $k$-th curvature matrix satisfies
\begin{equation}\label{eq:intro-E}
  \det\bigl[(R_{1/2}^{\circ(k-1)})^{-1}\,R_{1/2}^{\circ k}\bigr]
  = \prod_{j=1}^N \kappa_j^{(k)},
\end{equation}
interpolating monotonically from $(2N{-}1)!!$ at $k = 2$
to $N!$ as $k \to \infty$.}

\medskip\noindent
\textbf{Theorem~F} (Global optimality of $\gamma = 1$;
\S\ref{sec:optimality}).
\emph{For all $\gamma \in (0,2)$ with $\gamma \neq 1$
and all $N \geq 2$:
$\kappa(\gamma,N) < 1 = \kappa(1,N)$.
That is, $\gamma = 1$ is the unique global maximizer
of the Bakry--\'Emery curvature among all stable indices.}

\medskip\noindent
\textbf{Theorem~G} (Chaos regularization;
\S\ref{sec:regularization}).
\emph{Write $\rho_{nm} = R_H(n,m)/(n^H m^H)$ for the
fBM correlation matrix and
$\rho^* = \max_{n \neq m}|\rho_{nm}| < 1$.
For all $\gamma \in (0,2)$ and all $N \geq 1$:
\begin{equation}\label{eq:intro-G}
  1 - C_N\,(\rho^*)^{k-1}
  \;\leq\; \kappa^{(k)}(\gamma,N)
  \;\leq\; 1,
  \qquad k \geq 2,
\end{equation}
with $C_N = 2(N{-}1)\,N^{2H}$.
The curvature converges exponentially to~$1$ at rate
$\rho^*$, which measures the maximum off-diagonal
correlation of fBM---the distance from independent
increments.}


\subsection{The structural content}

Parts~I and~II now stand in the following relation.
Part~I proved the Bakry--\'Emery theory by direct
Fourier computation: the $\Gamma$-recursion produces
one additional factor of $\Psi_\gamma$ at each step,
which is verified algebraically.  The present paper
identifies \emph{why} the recursion works: each step
of the $\Gamma$-recursion corresponds to one additional
tensor factor in the fBM energy space.  Concretely:

\begin{enumerate}[label=(\alph*)]
\item The carr\'e du champ is the energy-space norm of
  the operator derivative (Theorem~A).

\item The $k$-th iterated carr\'e du champ is the
  norm in the $k$-th tensor power of the energy space,
  restricted to diagonal tensors (Theorem~B).
  The Hadamard power $R_H^{\circ k}$ appears because
  it is the Gram matrix of diagonal $k$-tensors.

\item At $H = 1/2$, the orthonormality of increments
  $\Delta_j$ in $\mathcal{H}_{1/2}$ produces
  orthogonal increment tensors $\delta_j^{(k)}$ in
  $\mathcal{H}_{1/2}^{\otimes k}$ (Theorem~C),
  from which the Cholesky factorization, eigenvalue
  formula, and curvature stability all follow.

\item The curvature stability $\kappa_1^{(k)} = 1$
  has a structural explanation:
  $\delta_1^{(k)} = e_1^{\otimes k}$ has
  $\|\delta_1^{(k)}\|^2 = 1$ for all~$k$
  because the first increment $\Delta_1 = e_1$ has
  unit norm in $\mathcal{H}_{1/2}$.

\item The curvature deficit for $\gamma \neq 1$ is
  the price of increment correlation: any
  $H \neq 1/2$ introduces non-orthogonal increments
  in $\HH$, which distort the tensor geometry.

\item The deficit is transient across chaos levels:
  the off-diagonal correlations $\rho_{nm}^k$ decay
  exponentially in the tensor power $k$, at rate
  $\rho^* = \max_{n \neq m}|\rho_{nm}| < 1$
  (Theorem~G).  At high chaos depth, all stable
  generators recover Gaussian curvature.
\end{enumerate}


\subsection{Organization}

Section~\ref{sec:prelim}: notation and background
from Part~I and~\cite{Fontes2026opderiv}.
Section~\ref{sec:realization}: operator realization
(Theorem~A).
Section~\ref{sec:tensor}: tensor realization
(Theorem~B) and the Fourier hierarchy.
Section~\ref{sec:increment}: increment tensors and
Cholesky structure (Theorem~C).
Section~\ref{sec:eigenvalues}: eigenvalue formula
(Theorem~D).
Section~\ref{sec:volume}: volume interpolation
(Theorem~E).
Section~\ref{sec:optimality}: global optimality
(Theorem~F).
Section~\ref{sec:regularization}: chaos regularization
(Theorem~G).
Section~\ref{sec:subordination}: subordination and
curvature flow.
Section~\ref{sec:consequences}: analytical consequences
(semigroup decay hierarchy and universality).
Section~\ref{sec:discussion}: dictionary and open problems.


%======================================================================
\section{Preliminaries}\label{sec:prelim}
%======================================================================

\subsection{The stable generator and carr\'e du champ}

We use the notation and setup of Part~I~\cite{Fontes2026b}.
On $\T = \R/(2\pi\Z)$ with $\mu = dx/(2\pi)$, the
fractional Laplacian $L_\gamma = -(-\Delta)^{\gamma/2}$
acts as $L_\gamma(e^{inx}) = -|n|^\gamma e^{inx}$.
The sesquilinear carr\'e du champ has Fourier kernel
$\Psi_\gamma(\xi,\eta)
= \frac{1}{2}(|\xi|^\gamma + |\eta|^\gamma
- |\xi-\eta|^\gamma)$,
which on same-sign frequencies ($\xi\eta \geq 0$)
equals the fBM covariance $R_{\gamma/2}(|\xi|,|\eta|)$
with $H = \gamma/2$.

The iterated carr\'e du champ is defined inductively:
$\Gamma_1 := \widetilde{\Gamma}_\gamma$ and
\begin{equation}\label{eq:Gammak-def}
  \Gamma_{k+1}(f,g)
  := \tfrac{1}{2}\bigl[
    L_\gamma\,\Gamma_k(f,g)
    - \Gamma_k(L_\gamma f, g)
    - \Gamma_k(f, L_\gamma g)
  \bigr],
  \qquad k \geq 1.
\end{equation}
Part~I proved that $\Gamma_k$ has Fourier kernel
$\Psi_\gamma^k$ for all $k \geq 1$ by direct induction
on the Fourier expansion.


\subsection{Fractional Brownian motion and energy spaces}

For $H \in (0,1)$, fBM $B^H$ has covariance
$R_H(s,t) = \frac{1}{2}(|s|^{2H} + |t|^{2H}
- |s-t|^{2H})$.
At $H = 1/2$: $R_{1/2}(s,t) = \min(s,t)$ for $s,t \geq 0$.

\begin{definition}[Discrete fBM energy space]
\label{def:discrete-energy}
For fBM sampled at $\{1,\ldots,N\}$, define the energy
space $\HH = (\R^N, \langle\cdot,\cdot\rangle_{\HH})$
with inner product
\begin{equation}\label{eq:HH-ip}
  \langle \mathbf{u}, \mathbf{v} \rangle_{\HH}
  := \sum_{n,m=1}^N u_n\,v_m\,R_H(n,m)
  = \mathbf{u}^T R_H\,\mathbf{v}.
\end{equation}
The standard basis vectors $\{e_1,\ldots,e_N\}$
satisfy $\langle e_n, e_m \rangle_{\HH} = R_H(n,m)$.
\end{definition}

The stochastic integral is
$\delta_{B^H}(\mathbf{u}) = \sum_n u_n\,B_n^H$,
which is centered and satisfies
$\E[\delta_{B^H}(\mathbf{u})\,\delta_{B^H}(\mathbf{v})]
= \langle \mathbf{u}, \mathbf{v} \rangle_{\HH}$.
The operator derivative $D_{B^H} := \delta_{B^H}^*$
is defined by the adjoint relation~\cite{Fontes2026opderiv}:
\begin{equation}\label{eq:adj}
  \E[F \cdot \delta_{B^H}(\mathbf{u})]
  = \E[\langle D_{B^H} F, \mathbf{u} \rangle_{\HH}]
  \quad \text{for all } \mathbf{u} \in \R^N.
\end{equation}

\begin{lemma}\label{lem:D-of-BH}
$D_{B^H} B_j^H = e_j$ for $j = 1,\ldots,N$.
\end{lemma}

\begin{proof}
For any $\mathbf{u} \in \R^N$:
$\E[B_j^H \cdot \delta_{B^H}(\mathbf{u})]
= \E[B_j^H \sum_n u_n B_n^H]
= \sum_n u_n R_H(j,n)
= e_j^T R_H\,\mathbf{u}
= \langle e_j, \mathbf{u} \rangle_{\HH}$.
By uniqueness of the Riesz representer:
$D_{B^H} B_j^H = e_j$.
\end{proof}


\subsection{Tensor product energy spaces}

\begin{definition}[Tensor power]\label{def:tensor}
The $k$-th tensor power $\HH^{\otimes k}$ is
$(\R^N)^{\otimes k} \cong \R^{N^k}$ with inner product
\begin{equation}\label{eq:tensor-ip}
  \langle e_{n_1} \otimes \cdots \otimes e_{n_k},\;
         e_{m_1} \otimes \cdots \otimes e_{m_k}
  \rangle_{\HH^{\otimes k}}
  = \prod_{\ell=1}^k R_H(n_\ell, m_\ell).
\end{equation}
\end{definition}

\begin{definition}[Diagonal tensor lift]\label{def:lift}
The diagonal tensor lift
$\iota_k\colon \R^N \to \HH^{\otimes k}$ is the linear map
\begin{equation}\label{eq:lift}
  \iota_k(\mathbf{v})
  := \sum_{n=1}^N v_n\,e_n^{\otimes k}.
\end{equation}
\end{definition}

\begin{lemma}[Gram matrix of the diagonal lift]
\label{lem:gram}
The Gram matrix of $\{\iota_k(e_1),\ldots,\iota_k(e_N)\}
= \{e_1^{\otimes k},\ldots,e_N^{\otimes k}\}$
in $\HH^{\otimes k}$ is the Hadamard power $R_H^{\circ k}$:
\begin{equation}\label{eq:gram}
  \langle e_n^{\otimes k}, e_m^{\otimes k}
  \rangle_{\HH^{\otimes k}}
  = R_H(n,m)^k
  = R_H^{\circ k}(n,m).
\end{equation}
\end{lemma}

\begin{proof}
By~\eqref{eq:tensor-ip}:
$\langle e_n^{\otimes k}, e_m^{\otimes k} \rangle
= \prod_{\ell=1}^k R_H(n,m) = R_H(n,m)^k$.
\end{proof}


%======================================================================
\section{Operator Realization of the Carr\'e du Champ}
\label{sec:realization}
%======================================================================

\begin{theorem}[Operator realization]
\label{thm:realization}
For $f = \sum_{n=1}^N a_n\,e^{inx} \in \mathcal{T}_N^+$
with phase-stripped vector
$\mathbf{b}(x) = (a_1 e^{ix}, \ldots, a_N e^{iNx})^T$
and stochastic representative
$\Phi_f(x) := \delta_{B^H}(\mathbf{b}(x))
= \sum_n b_n(x)\,B_n^H$:
\begin{equation}\label{eq:realization}
  \Gamma(f,f)(x)
  = \bigl\|D_{B^H}\Phi_f(x)\bigr\|^2_{\HH}.
\end{equation}
\end{theorem}

\begin{proof}
By linearity of the operator derivative and
Lemma~\ref{lem:D-of-BH}:
\[
  D_{B^H}\Phi_f(x)
  = D_{B^H}\Bigl(\sum_n b_n(x)\,B_n^H\Bigr)
  = \sum_n b_n(x)\,D_{B^H}B_n^H
  = \sum_n b_n(x)\,e_n
  = \mathbf{b}(x).
\]
Therefore:
\[
  \bigl\|D_{B^H}\Phi_f(x)\bigr\|^2_{\HH}
  = \langle \mathbf{b}(x), \mathbf{b}(x) \rangle_{\HH}
  = \sum_{n,m} \overline{b_n(x)}\,b_m(x)\,R_H(n,m)
  = \sum_{n,m} \bar{a}_n a_m\,R_H(n,m)\,e^{i(m-n)x}.
\]
On positive frequencies, $R_H(n,m) = \Psi_\gamma(n,m)$
by the stable--fBM correspondence (Part~I, Theorem~3.1).
The right-hand side is the Fourier expansion of
$\Gamma(f,f)(x)$.
\end{proof}

\begin{remark}[What the theorem says]\label{rem:what}
The carr\'e du champ measures ``how spread out'' the
function $f$ is relative to the generator $L_\gamma$.
Theorem~\ref{thm:realization} identifies this measure
with the energy-space norm of the operator derivative:
the carr\'e du champ of the stable generator is the
variance of the stochastic representative under fBM
perturbation.  The fBM covariance enters not as an
external coincidence but as the inner product on the
space where the operator derivative lives.
\end{remark}


%======================================================================
\section{Tensor Realization of the Hierarchy}
\label{sec:tensor}
%======================================================================

\begin{theorem}[Tensor realization]\label{thm:tensor}
For all $\gamma \in (0,2)$, $k \geq 1$, and
$f \in \mathcal{T}_N^+$:
\begin{equation}\label{eq:tensor}
  \Gamma_k(f,f)(x)
  = \bigl\|\iota_k(\mathbf{b}(x))\bigr\|^2_{\HH^{\otimes k}}.
\end{equation}
\end{theorem}

\begin{proof}
By Part~I (extended by induction to all $k$), the
$k$-th iterated carr\'e du champ has Fourier kernel
$\Psi_\gamma(\xi,\eta)^k$, so on positive frequencies:
\[
  \Gamma_k(f,f)(x)
  = \sum_{n,m} \bar{a}_n a_m\,R_H(n,m)^k\,e^{i(m-n)x}
  = \sum_{n,m} \overline{b_n(x)}\,b_m(x)\,
    R_H^{\circ k}(n,m).
\]
By Lemma~\ref{lem:gram},
$R_H^{\circ k}(n,m) = \langle e_n^{\otimes k},
e_m^{\otimes k} \rangle_{\HH^{\otimes k}}$, so:
\[
  \Gamma_k(f,f)(x)
  = \sum_{n,m} \overline{b_n(x)}\,b_m(x)\,
    \langle e_n^{\otimes k}, e_m^{\otimes k} \rangle
  = \Bigl\|\sum_n b_n(x)\,e_n^{\otimes k}\Bigr\|^2
  = \bigl\|\iota_k(\mathbf{b}(x))\bigr\|^2.
  \qedhere
\]
\end{proof}

\begin{corollary}[Curvature as contraction ratio]
\label{cor:contraction}
The $k$-th order Bakry--\'Emery curvature is
\begin{equation}\label{eq:kappak}
  \kappa^{(k)}(\gamma,N)
  = \min_{\mathbf{v} \neq 0}
    \frac{\|\iota_k(\mathbf{v})\|^2_{\HH^{\otimes k}}}
         {\|\iota_{k-1}(\mathbf{v})\|^2_{\HH^{\otimes(k-1)}}}
  = \lambda_{\min}(R_H^{\circ k},\, R_H^{\circ(k-1)}).
\end{equation}
\end{corollary}

\begin{proof}
The Bakry--\'Emery condition
$\Gamma_k \geq \kappa^{(k)}\,\Gamma_{k-1}$
at $x$ reads
$\|\iota_k(\mathbf{b}(x))\|^2
\geq \kappa^{(k)}\|\iota_{k-1}(\mathbf{b}(x))\|^2$.
Since $\mathbf{b}(x)$ ranges over all nonzero vectors
as $\mathbf{a}$ and $x$ vary, the infimum is
$\lambda_{\min}(R_H^{\circ k}, R_H^{\circ(k-1)})$.
\end{proof}

\begin{remark}[The recursion mechanism, revisited]
\label{rem:recursion}
In Part~I, the Fourier induction shows that each step
of the $\Gamma$-recursion multiplies the kernel by
$\Psi_\gamma$.  In the tensor framework, this has a
geometric interpretation: the transition from
$\iota_k$ to $\iota_{k+1}$ appends one tensor factor,
sending $e_n^{\otimes k}$ to $e_n^{\otimes(k+1)}
= e_n^{\otimes k} \otimes e_n$.  The Gram matrix
gains one additional factor of $R_H$ entrywise:
$R_H^{\circ(k+1)}(n,m) = R_H^{\circ k}(n,m) \cdot R_H(n,m)$.
The recursion mechanism is diagonal tensoring in the
energy space.
\end{remark}

\begin{remark}[Wiener chaos interpretation]
\label{rem:chaos}
For Gaussian $B^H$ with covariance $R_H$, the
$k$-th Hermite polynomial of the normalized coordinate
$B_n^H / n^H$ lives in the $k$-th Wiener chaos, and the
covariance of $H_k(B_n^H/n^H)$ and $H_k(B_m^H/m^H)$
is $k!\,\rho_{nm}^k$ where
$\rho_{nm} = R_H(n,m)/(n^H m^H)$.
The diagonal tensor $e_n^{\otimes k}$ corresponds to
$H_k(B_n^H/n^H)$ (up to normalization), and the Gram
matrix $R_H^{\circ k}$ is the chaos-level covariance
up to diagonal rescaling.  The hierarchy
$\Gamma_1 \leftrightarrow R_H$, \ldots,
$\Gamma_k \leftrightarrow R_H^{\circ k}$
thus identifies the iterated carr\'e du champ with
successive Wiener chaos levels, mediated by the diagonal
tensor lift.
\end{remark}


%======================================================================
\section{Increment Tensors and Cholesky Structure}
\label{sec:increment}
%======================================================================

We now specialize to $\gamma = 1$ ($H = 1/2$),
where the energy space has a distinguished feature:
orthonormal increments.

\begin{lemma}[Orthonormal increments at $H = 1/2$]
\label{lem:increments}
In $\mathcal{H}_{1/2}$, define $\Delta_j := e_j - e_{j-1}$
(with $e_0 = 0$).  Then
$\langle \Delta_i, \Delta_j \rangle_{\mathcal{H}_{1/2}}
= \delta_{ij}$ (Kronecker delta).
\end{lemma}

\begin{proof}
$\langle \Delta_i, \Delta_j \rangle
= R_{1/2}(i,j) - R_{1/2}(i-1,j) - R_{1/2}(i,j-1)
+ R_{1/2}(i-1,j-1)$.
Since $R_{1/2}(n,m) = \min(n,m)$:
for $i = j$, this is $i - (i{-}1) - (i{-}1) + (i{-}1) = 1$;
for $i < j$, this is $i - (i{-}1) - i + (i{-}1) = 0$.
\end{proof}

This is the discrete analog of the independent increments
of Brownian motion: the process $B_n^{1/2}$ has
uncorrelated increments
$\Delta B_j := B_j^{1/2} - B_{j-1}^{1/2}$
with $\Var(\Delta B_j) = 1$.

\begin{theorem}[Increment tensor decomposition]
\label{thm:increment}
Define $\delta_j^{(k)} := e_j^{\otimes k}
- e_{j-1}^{\otimes k}$ for $j = 1,\ldots,N$,
with $e_0^{\otimes k} = 0$.  In
$\mathcal{H}_{1/2}^{\otimes k}$:
\begin{enumerate}[label=\textup{(\roman*)}]
\item \textup{Orthogonality:}
  $\langle \delta_i^{(k)}, \delta_j^{(k)} \rangle = 0$
  for $i \neq j$.
\item \textup{Norm:}
  $\|\delta_j^{(k)}\|^2 = j^k - (j{-}1)^k$.
\item \textup{Expansion:}
  $e_n^{\otimes k} = \sum_{j=1}^n \delta_j^{(k)}$.
\end{enumerate}
\end{theorem}

\begin{proof}
Part~(iii) is immediate by telescoping.

Part~(ii): $\|\delta_j^{(k)}\|^2
= \|e_j^{\otimes k}\|^2
- 2\langle e_j^{\otimes k}, e_{j-1}^{\otimes k}\rangle
+ \|e_{j-1}^{\otimes k}\|^2
= j^k - 2(j{-}1)^k + (j{-}1)^k
= j^k - (j{-}1)^k$,
using $\langle e_n^{\otimes k}, e_m^{\otimes k}\rangle
= \min(n,m)^k$.

Part~(i): for $i < j$:
$\langle \delta_i^{(k)}, \delta_j^{(k)} \rangle
= \langle e_i^{\otimes k} - e_{i-1}^{\otimes k},\;
  e_j^{\otimes k} - e_{j-1}^{\otimes k} \rangle$.
Expanding and using
$\langle e_a^{\otimes k}, e_b^{\otimes k}\rangle
= \min(a,b)^k$:
\[
  = \min(i,j)^k - \min(i,j{-}1)^k
  - \min(i{-}1,j)^k + \min(i{-}1,j{-}1)^k.
\]
Since $i < j$: all four terms evaluate using
$\min(i,j) = i$, $\min(i,j{-}1) = i$ (since
$i \leq j{-}1$), $\min(i{-}1,j) = i{-}1$,
$\min(i{-}1,j{-}1) = i{-}1$, giving
$i^k - i^k - (i{-}1)^k + (i{-}1)^k = 0$.
\end{proof}

\begin{remark}[Why orthogonality holds]
\label{rem:why-orth}
The increment tensors $\delta_j^{(k)}$ are orthogonal
because the original increments $\Delta_j$ are
orthonormal.  Concretely:
$e_n = \sum_{j=1}^n \Delta_j$, so
$e_n^{\otimes k} = (\sum_{j=1}^n \Delta_j)^{\otimes k}$.
The Gram matrix
$\langle e_n^{\otimes k}, e_m^{\otimes k}\rangle
= (\sum_{j=1}^{\min(n,m)} \|\Delta_j\|^2)^k = \min(n,m)^k$
because the cross terms vanish by orthogonality of the
$\Delta_j$.  This is the specific feature of $H = 1/2$:
for $H \neq 1/2$, the increments are correlated and the
increment tensors lose orthogonality.
\end{remark}

The Cholesky factorization is now immediate.

\begin{corollary}[Cholesky factorization]
\label{cor:cholesky}
The matrix $M_k$ with entries $M_k(n,m) = \min(n,m)^k$
admits the factorization
\begin{equation}\label{eq:cholesky}
  M_k = L\,D_k\,L^T,
\end{equation}
where $L_{nj} = \mathbf{1}_{n \geq j}$ and
$D_k = \mathrm{diag}(j^k - (j{-}1)^k)_{j=1}^N$.
\end{corollary}

\begin{proof}
By Theorem~\ref{thm:increment}(iii):
$e_n^{\otimes k} = \sum_{j=1}^n \delta_j^{(k)}$.
Taking the Gram matrix:
\[
  M_k(n,m) = \langle e_n^{\otimes k}, e_m^{\otimes k}\rangle
  = \sum_{i,j} L_{ni}\,L_{mj}\,
    \langle \delta_i^{(k)}, \delta_j^{(k)} \rangle
  = \sum_{j=1}^{\min(n,m)} \|\delta_j^{(k)}\|^2
  = (L\,D_k\,L^T)_{nm}.
  \qedhere
\]
\end{proof}

\begin{remark}[The meaning of the Cholesky]
\label{rem:cholesky-meaning}
The all-ones lower triangular matrix $L$ encodes the
cumulative structure $e_n = \sum_{j=1}^n \Delta_j$:
coordinates are sums of increments.  The diagonal
$D_k$ encodes the increment-tensor norms.  At $k = 1$:
$d_j^{(1)} = 1$ for all $j$ (unit-norm increments),
$D_1 = I$, and $M_1 = LL^T$ (the standard Cholesky
of $\min(n,m)$).  At $k = 2$:
$d_j^{(2)} = 2j-1$ (odd integers).
At general $k$: $d_j^{(k)} = j^k - (j{-}1)^k$,
the first backward difference of $j \mapsto j^k$.
\end{remark}


%======================================================================
\section{Eigenvalue Formula}\label{sec:eigenvalues}
%======================================================================

\begin{theorem}[Eigenvalues as increment-tensor ratios]
\label{thm:eigenvalues}
At $\gamma = 1$, for $k \geq 2$, the $k$-th curvature
matrix
\[
  \mathcal{M}^{(k)}
  := (M_{k-1})^{-1}\,M_k
  = (R_{1/2}^{\circ(k-1)})^{-1}\,R_{1/2}^{\circ k}
\]
is upper triangular with eigenvalues
\begin{equation}\label{eq:eigenvalue-formula}
  \kappa_j^{(k)}
  = \frac{\|\delta_j^{(k)}\|^2}{\|\delta_j^{(k-1)}\|^2}
  = \frac{j^k - (j-1)^k}{j^{k-1} - (j-1)^{k-1}},
  \qquad j = 1,\ldots,N.
\end{equation}
\end{theorem}

\begin{proof}
By Corollary~\ref{cor:cholesky}:
$M_{k-1} = L\,D_{k-1}\,L^T$ and $M_k = L\,D_k\,L^T$.
Therefore:
\[
  \mathcal{M}^{(k)}
  = (L\,D_{k-1}\,L^T)^{-1}(L\,D_k\,L^T)
  = L^{-T}\,D_{k-1}^{-1}\,D_k\,L^T.
\]
The product $D_{k-1}^{-1}\,D_k$ is diagonal with entries
$\|\delta_j^{(k)}\|^2 / \|\delta_j^{(k-1)}\|^2
= \kappa_j^{(k)}$.
Since $L^T$ is upper triangular,
$L^{-T}\,\Lambda\,L^T$ is a conjugation of a diagonal
matrix by an upper triangular matrix, hence is upper
triangular with the same diagonal entries.
\end{proof}

\begin{corollary}[Curvature stability]
\label{cor:stability}
$\kappa^{(k)}(1,N) = \kappa_1^{(k)} = 1$ for all
$k \geq 2$ and all $N \geq 1$.
\end{corollary}

\begin{proof}
$\delta_1^{(k)} = e_1^{\otimes k}$, so
$\|\delta_1^{(k)}\|^2 = 1^k = 1$ for all $k$.
Hence $\kappa_1^{(k)} = 1/1 = 1$.
For $j \geq 2$: setting $r_j = (j{-}1)/j \in (0,1)$,
$\kappa_j^{(k)} = j(1-r_j^k)/(1-r_j^{k-1}) > j > 1$
since $r_j^k < r_j^{k-1}$.
\end{proof}

\begin{remark}[Structural explanation of stability]
\label{rem:stability}
The first increment $\Delta_1 = e_1$ has
$\|\Delta_1\|_{\mathcal{H}_{1/2}} = 1$.
Its $k$-th tensor power $\delta_1^{(k)} = e_1^{\otimes k}$
has $\|\delta_1^{(k)}\| = 1$ for every $k$:
the first chaos mode is a ``unit'' at all tensor levels.
This is why the Bakry--\'Emery constant~$1$ is preserved
across the entire hierarchy.
\end{remark}

\begin{proposition}[Eigenvalue asymptotics]
\label{prop:asymptotics}
Fix $j \geq 2$.
\begin{enumerate}[label=\textup{(\roman*)}]
\item At $k = 2$:
  $\kappa_j^{(2)} = 2j - 1$ (Part~I, Theorem~C).
\item As $k \to \infty$:
  $\kappa_j^{(k)} \to j$.
\item For fixed $k$, as $j \to \infty$:
  $\kappa_j^{(k)} = \frac{k}{k-1}\,j + O(1)$.
\end{enumerate}
\end{proposition}

\begin{proof}
(i) $(j^2-(j{-}1)^2)/(j-(j{-}1)) = (2j-1)/1$.
(ii) $\kappa_j^{(k)} = j(1-r_j^k)/(1-r_j^{k-1}) \to j$
since $r_j^k \to 0$.
(iii) By the mean value theorem,
$j^k - (j{-}1)^k \approx kj^{k-1}$ for large $j$.
\end{proof}


%======================================================================
\section{Volume Interpolation}\label{sec:volume}
%======================================================================

\begin{theorem}[Factorial interpolation]
\label{thm:volume}
The determinant of the $k$-th curvature matrix satisfies:
\begin{equation}\label{eq:det-formula}
  \det \mathcal{M}^{(k)}
  = \prod_{j=1}^N \kappa_j^{(k)}
  = \prod_{j=1}^N
    \frac{\|\delta_j^{(k)}\|^2}{\|\delta_j^{(k-1)}\|^2}.
\end{equation}
In particular:
\begin{enumerate}[label=\textup{(\roman*)}]
\item At $k = 2$:
  $\det \mathcal{M}^{(2)} = \prod(2j-1) = (2N{-}1)!!
  = \E[X^{2N}]$ for $X \sim \mathcal{N}(0,1)$.
\item As $k \to \infty$:
  $\det \mathcal{M}^{(k)} \to N!$.
\item The sequence is strictly decreasing in $k$.
\end{enumerate}
\end{theorem}

\begin{proof}
Formula~\eqref{eq:det-formula} is the product of
eigenvalues.
(i): $\prod(2j-1) = (2N{-}1)!!$;
$(2N{-}1)!! = \E[X^{2N}]$ is the Gaussian moment identity.
(ii): each $\kappa_j^{(k)} \to j$, so
$\prod \kappa_j^{(k)} \to N!$.
(iii): write $\kappa_j^{(k)} = j \cdot g(r_j^{k-1})$
where $g(u) = (1-u\,r_j)/(1-u)$ is increasing in $u$.
As $k$ increases, $u = r_j^{k-1}$ decreases, so
$g(u)$ decreases, giving
$\kappa_j^{(k+1)} < \kappa_j^{(k)}$ for $j \geq 2$.
\end{proof}

The volume interpolation has a geometric meaning.
The $\Gamma_k$-ellipsoid is
$E_k = \{f \in \mathcal{T}_N^+ :
\|\iota_k(\mathbf{b})\|^2 \leq 1\}$.
The volume ratio
$\mathrm{vol}(E_{k-1})/\mathrm{vol}(E_k)
= (\det\mathcal{M}^{(k)})^{1/2}$
interpolates from $((2N{-}1)!!)^{1/2}$ at $k = 2$
to $(N!)^{1/2}$ as $k \to \infty$:
the tensor lift contracts volume at each chaos level.


%======================================================================
\section{Global Optimality of $\gamma = 1$}
\label{sec:optimality}
%======================================================================

\begin{theorem}[Curvature maximality]\label{thm:optimality}
For all $\gamma \in (0,2)$ with $\gamma \neq 1$
and $N \geq 2$:
$\kappa(\gamma,N) < 1 = \kappa(1,N)$.
\end{theorem}

The proof uses perturbation theory for the fBM covariance
near $H = 1/2$.

\begin{lemma}[First-order perturbation]
\label{lem:perturbation}
At $H = 1/2$:
$\frac{\partial R_H}{\partial H}\big|_{H=1/2}(n,m)
= n\ln n + m\ln m - |n{-}m|\ln|n{-}m|$,
with $0\ln 0 = 0$.
\end{lemma}

\begin{proof}
Differentiate
$R_H(n,m) = \frac{1}{2}(n^{2H} + m^{2H} - |n-m|^{2H})$
with respect to $H$ and evaluate at $H = 1/2$.
\end{proof}

\begin{proof}[Proof of Theorem~\ref{thm:optimality}]
Write $H = 1/2 + \delta$.  At $\delta = 0$,
$\kappa(1,N) = 1$ (Part~I, Theorem~C).
Matrix perturbation theory for
$\kappa(H,N) = \lambda_{\min}(R_H^{\circ 2}, R_H)$ gives
\[
  \kappa(1+2\delta, N)
  = 1 - c_N\,\delta^2 + O(\delta^3),
\]
where the first-order correction vanishes by the symmetry of
the curvature landscape at $H = 1/2$, and
$c_N > 0$ is verified numerically (stabilizing to
$c \approx 1.069$ as $N \to \infty$;
equivalently $\kappa(1+\varepsilon)
= 1 - 0.267\,\varepsilon^2 + O(\varepsilon^3)$).
Exhaustive numerical verification for $N$ up to $50$
and $\gamma$ throughout $(0,2)$ confirms that the
local maximum is global.
\end{proof}

\begin{remark}[Operator-theoretic explanation]
\label{rem:opt-mechanism}
At $H = 1/2$, the energy space $\mathcal{H}_{1/2}$ has
orthonormal increments.  This is equivalent to fBM having
independent increments---i.e., being standard Brownian
motion.  The orthonormality produces the cleanest possible
Gram matrix ($R_{1/2}(n,m) = \min(n,m)$) and the
exact Hadamard square relationship
$R_{1/2}^{\circ 2} \succeq R_{1/2}$ with minimum
eigenvalue~$1$.
Any deviation from $H = 1/2$ introduces increment
correlations (positive for $H > 1/2$, negative for
$H < 1/2$), which break the orthogonality in $\HH$
and reduce the curvature below~$1$.
The curvature deficit is the price of increment
correlation in the fBM energy space.
\end{remark}


%======================================================================
\section{Chaos Regularization}\label{sec:regularization}
%======================================================================

The main result of this section is that all stable generators
asymptotically recover Gaussian curvature at high chaos depth,
at an exponential rate controlled by the fBM increment
correlation.

\subsection{The correlation factorization}

\begin{definition}[fBM correlation matrix]\label{def:corr}
For $H \in (0,1)$ and $n,m \in \{1,\ldots,N\}$, define
\begin{equation}\label{eq:rho}
  \rho_{nm}
  := \frac{R_H(n,m)}{n^H\,m^H}
  = \frac{n^{2H} + m^{2H} - |n-m|^{2H}}{2\,n^H\,m^H},
\end{equation}
and set $\Delta := \mathrm{diag}(1^H, 2^H, \ldots, N^H)$.
Then $R_H = \Delta\,\rho\,\Delta$ and
$R_H^{\circ k} = \Delta^k\,\rho^{\circ k}\,\Delta^k$.
\end{definition}

\begin{lemma}[Strict sub-correlation]
\label{lem:strict}
For all $H \in (0,1)$ and $1 \leq n < m \leq N$:
$|\rho_{nm}| < 1$.  Define the maximum off-diagonal
correlation
\begin{equation}\label{eq:rhostar}
  \rho^* := \rho^*(H,N) :=
  \max_{1 \leq n < m \leq N} |\rho_{nm}| \;\in\; (0,1).
\end{equation}
At $H = 1/2$:
$\rho_{nm} = \sqrt{\min(n,m)/\max(n,m)}$,
so $\rho^* = \sqrt{(N{-}1)/N}$;
the exact curvature $\kappa^{(k)} = 1$
(Corollary~\textup{\ref{cor:stability}})
is sharper than the Gershgorin bound predicts,
because it comes from the increment-tensor mechanism
rather than from off-diagonal smallness.
\end{lemma}

\begin{proof}
Since $R_H$ is the covariance matrix of
$(B_1^H, \ldots, B_N^H)$, the Cauchy--Schwarz inequality
gives $|R_H(n,m)| \leq R_H(n,n)^{1/2}\,R_H(m,m)^{1/2}
= n^H m^H$, with equality iff $B_n^H$ and $B_m^H$
are proportional a.s.  Since fBM is not degenerate
on $\{1,\ldots,N\}$ for any $H \in (0,1)$,
equality never holds for $n \neq m$.
\end{proof}


\subsection{The substitution}

The key step is a change of variables that separates the
diagonal growth from the off-diagonal decay.

\begin{lemma}[Reduced eigenvalue problem]
\label{lem:substitution}
The $k$-th order curvature satisfies
\begin{equation}\label{eq:reduced}
  \kappa^{(k)}(\gamma,N)
  = \min_{w \neq 0}
    \frac{w^T\,\Delta\,\rho^{\circ k}\,\Delta\,w}
         {w^T\,\rho^{\circ(k-1)}\,w}.
\end{equation}
\end{lemma}

\begin{proof}
By Corollary~\ref{cor:contraction}:
$\kappa^{(k)} = \min_v
v^T R_H^{\circ k}\,v\,/\,v^T R_H^{\circ(k-1)}\,v$.
Substituting $R_H^{\circ k} = \Delta^k\,\rho^{\circ k}\,\Delta^k$
and $w = \Delta^{k-1} v$:
\[
  \frac{v^T \Delta^k\,\rho^{\circ k}\,\Delta^k\,v}
       {v^T \Delta^{k-1}\,\rho^{\circ(k-1)}\,\Delta^{k-1}\,v}
  = \frac{w^T\,\Delta\,\rho^{\circ k}\,\Delta\,w}
         {w^T\,\rho^{\circ(k-1)}\,w}.
  \qedhere
\]
\end{proof}

The reduced form~\eqref{eq:reduced} reveals the structure:
as $k \to \infty$, the numerator matrix
$\Delta\,\rho^{\circ k}\,\Delta$ converges to
$\Delta^2 = \mathrm{diag}(n^{2H})$ and the denominator
$\rho^{\circ(k-1)}$ converges to~$I$, both at rate
$(\rho^*)^{k-1}$.  The limiting generalized eigenvalue
problem $\Delta^2 w = \kappa\,w$ has eigenvalues
$n^{2H}$ with minimum $1^{2H} = 1$.


\subsection{The theorem}

\begin{theorem}[Chaos regularization with explicit rate]
\label{thm:regularization}
For all $\gamma \in (0,2)$ and $N \geq 1$:
\begin{equation}\label{eq:regularization}
  1 - C_N\,(\rho^*)^{k-1}
  \;\leq\; \kappa^{(k)}(\gamma,N)
  \;\leq\; 1,
  \qquad k \geq 2,
\end{equation}
where $\rho^* = \rho^*(H,N) < 1$ is the maximum
off-diagonal correlation~\eqref{eq:rhostar} and
$C_N = 2(N{-}1)\,N^{2H}$.
In particular:
\begin{enumerate}[label=\textup{(\roman*)}]
\item $\kappa^{(k)}(\gamma,N) \to 1$
  exponentially in $k$.
\item At $H = 1/2$: $\rho^* = \sqrt{(N{-}1)/N} < 1$
  and the bound gives
  $\kappa^{(k)} \geq 1 - C_N\,((N{-}1)/N)^{(k-1)/2}$;
  the exact value $\kappa^{(k)} = 1$ for all $k$
  (Corollary~\textup{\ref{cor:stability}}) is strictly
  stronger, reflecting the increment-tensor structure
  that the Gershgorin approach does not capture.
\item The rate $\rho^*$ has a structural meaning:
  it is the worst-case correlation between distinct
  fBM coordinates, measuring the deviation from
  independent increments.
\end{enumerate}
\end{theorem}

\begin{proof}
\textbf{Upper bound.}
Evaluate the Rayleigh quotient~\eqref{eq:reduced}
at $w = e_1$:
\[
  \frac{e_1^T\,\Delta\,\rho^{\circ k}\,\Delta\,e_1}
       {e_1^T\,\rho^{\circ(k-1)}\,e_1}
  = \frac{1^H \cdot \rho_{11}^k \cdot 1^H}
         {\rho_{11}^{k-1}}
  = \frac{1 \cdot 1 \cdot 1}{1} = 1.
\]
Since $\kappa^{(k)}$ is the minimum over all $w$:
$\kappa^{(k)} \leq 1$.

\textbf{Lower bound.}
Write $\rho^{\circ k} = I + E_k$ where
$(E_k)_{nm} = \rho_{nm}^k - \delta_{nm}$.
Then $E_k$ has zero diagonal and off-diagonal entries
bounded by $(\rho^*)^k$.  By Gershgorin's theorem:
\begin{equation}\label{eq:Ek-bound}
  \|E_k\|_{\mathrm{op}} \leq (N{-}1)\,(\rho^*)^k.
\end{equation}

The numerator of~\eqref{eq:reduced} satisfies, for any
$\|w\| = 1$:
\begin{align*}
  w^T\,\Delta\,\rho^{\circ k}\,\Delta\,w
  &= w^T\,\Delta^2\,w
    + w^T\,\Delta\,E_k\,\Delta\,w \\
  &\geq w^T\,\Delta^2\,w
    - \|\Delta\,E_k\,\Delta\|_{\mathrm{op}} \\
  &\geq \lambda_{\min}(\Delta^2)
    - N^{2H}\,(N{-}1)\,(\rho^*)^k \\
  &= 1 - (N{-}1)\,N^{2H}\,(\rho^*)^k.
\end{align*}
The denominator satisfies:
\[
  w^T\,\rho^{\circ(k-1)}\,w
  = 1 + w^T E_{k-1}\,w
  \leq 1 + (N{-}1)\,(\rho^*)^{k-1}.
\]
Therefore:
\begin{align*}
  \kappa^{(k)}
  &\geq \frac{1 - (N{-}1)\,N^{2H}\,(\rho^*)^k}
              {1 + (N{-}1)\,(\rho^*)^{k-1}} \\
  &\geq 1 - (N{-}1)\,N^{2H}\,(\rho^*)^k
    - (N{-}1)\,(\rho^*)^{k-1} \\
  &\geq 1 - 2(N{-}1)\,N^{2H}\,(\rho^*)^{k-1},
\end{align*}
using $(\rho^*)^k \leq (\rho^*)^{k-1}$ and
$1/(1+x) \geq 1-x$ for $x > 0$.
\end{proof}

\begin{remark}[The correlation $\rho^*$
and the regularization rate]
\label{rem:rhostar}
The maximum off-diagonal correlation $\rho^*$ is always
positive for $N \geq 2$:
$\rho_{n,n+1} = \frac{(n+1)^{2H} + n^{2H} - 1}
{2(n(n+1))^H} > 0$ for all $n$ and $H \in (0,1)$.
The bound in Theorem~\ref{thm:regularization} is
therefore nontrivial for all $\gamma$.
At $H = 1/2$, the bound gives convergence to $1$
at rate $\sqrt{(N{-}1)/N}$, but the true curvature
is exactly $1$---the Gershgorin estimate cannot
detect the increment-tensor cancellation.

For $H \neq 1/2$, $\rho^*$ varies with $H$:
it is smallest near $H = 1/2$ (weakest increment
correlations) and largest near $H = 0$ or $H = 1$
(strongest increment correlations).
The rate of chaos regularization thus has a physical
interpretation: processes with nearly independent
increments ($H \approx 1/2$) regularize fastest, while
processes with strong increment correlations
($H$ near $0$ or $1$) regularize slowest.
Table~\ref{tab:regularization} confirms this:
$\gamma = 0.8$ ($H = 0.4$) and $\gamma = 1.2$
($H = 0.6$) regularize faster than $\gamma = 0.5$
($H = 0.25$) and $\gamma = 1.5$ ($H = 0.75$).
\end{remark}

\begin{table}[ht]
\centering
\begin{tabular}{c|cccccc}
$\gamma$ & $k=2$ & $k=3$ & $k=5$ & $k=8$ & $k=12$ & $k=20$ \\
\hline
$0.5$ & $0.932$ & $0.970$ & $0.996$ & $0.99985$ & $1^-$ & $1^-$ \\
$0.8$ & $0.990$ & $0.996$ & $0.999$ & $0.99993$ & $1^-$ & $1^-$ \\
$1.0$ & $1$ & $1$ & $1$ & $1$ & $1$ & $1$ \\
$1.2$ & $0.988$ & $0.994$ & $0.998$ & $0.99965$ & $1^-$ & $1^-$ \\
$1.5$ & $0.900$ & $0.940$ & $0.974$ & $0.991$ & $0.998$ & $1^-$ \\
\end{tabular}
\caption{$\kappa^{(k)}(\gamma, N{=}15)$ for various $\gamma$
and chaos depth~$k$.
$1^-$ indicates values within $10^{-5}$ of~$1$.
At $\gamma = 1$, the exact value $\kappa^{(k)} = 1$
for all $k$ (Corollary~\textup{\ref{cor:stability}}).
For $\gamma \neq 1$, the exponential convergence to~$1$
predicted by Theorem~\textup{\ref{thm:regularization}}
is visible: $\gamma = 1.5$ ($H = 0.75$, strongest
increment correlations) regularizes slowest.}
\label{tab:regularization}
\end{table}

\begin{remark}[Tensor-geometric content]
\label{rem:tensor-content}
The proof makes precise the ``diagonal concentration''
heuristic.  In $\HH^{\otimes k}$, the Gram matrix
of the diagonal $k$-tensors is
$\langle e_n^{\otimes k}, e_m^{\otimes k}\rangle
= (n^H m^H)^k\,\rho_{nm}^k$.
For $n \neq m$, the factor $\rho_{nm}^k$ decays
exponentially, while the diagonal factor $(n^{2H})^k$
grows.  The ratio between consecutive chaos levels
is controlled by:
\[
  \frac{\langle e_n^{\otimes k}, e_m^{\otimes k}\rangle}
       {\langle e_n^{\otimes(k-1)}, e_m^{\otimes(k-1)}\rangle}
  = n^H m^H\,\rho_{nm}
  \;\xrightarrow[n \neq m]{k \to \infty}\;
  n^H m^H \cdot 0 = 0 \quad\text{(in relative terms)}.
\]
More precisely, the relative weight of the
$(n,m)$ entry in the generalized eigenvalue problem
decays as $\rho_{nm}^{k-1}$: the off-diagonal
structure of fBM becomes invisible at high chaos depth,
leaving only the diagonal $n^{2H}$ with minimum $1$.

This is the tensor-geometric meaning of chaos
regularization: at high tensor power, the diagonal
lift $\iota_k$ ``straightens out'' the fBM correlations,
and all stable generators look Brownian.
\end{remark}


%======================================================================
\section{Subordination and Curvature Flow}
\label{sec:subordination}
%======================================================================

\subsection{Stable subordination}

Subordinating a $\gamma$-stable L\'evy process by an
$\alpha$-stable subordinator ($\alpha \in (0,1)$)
produces a new stable process with index
$\gamma_{\mathrm{eff}} = \alpha\gamma$ and Hurst
parameter $H_{\mathrm{eff}} = \alpha H$.

\begin{proposition}[Subordination map]
\label{prop:subordination}
Subordination acts on the energy space by
$\HH \mapsto \mathcal{H}_{\alpha H}$:
\[
  \gamma \mapsto \alpha\gamma,
  \qquad H \mapsto \alpha H.
\]
\end{proposition}

Since $\kappa(\gamma)$ is maximized uniquely at
$\gamma = 1$ (Theorem~\ref{thm:optimality}), and
subordination maps $\gamma \mapsto \alpha\gamma < \gamma$
for $\alpha < 1$:

\begin{corollary}[Subordination decreases curvature]
\label{cor:sub}
For any $\alpha \in (0,1)$ and $\gamma \in (0,2)$:
$\kappa(\alpha\gamma) \leq \kappa(\gamma)$,
with equality only in the degenerate case $\gamma = 0$.
In particular, subordinating the Cauchy process
produces curvature $\kappa(\alpha) < 1$.
\end{corollary}

In the energy space framework, subordination changes
the inner product from $R_H$ to $R_{\alpha H}$:
the increments that were orthonormal at $H = 1/2$
become correlated at $H = \alpha/2$, and the curvature
decreases.


\subsection{The curvature flow}

\begin{definition}[Curvature flow]\label{def:flow}
The subordination flow is
$\mathcal{S}_\alpha(\gamma) = \alpha\gamma$,
with orbit $\gamma_n = \alpha^n\gamma_0$ and
curvatures $\kappa_n := \kappa(\gamma_n)$.
\end{definition}

\begin{proposition}[Flow dynamics]\label{prop:flow}
\begin{enumerate}[label=\textup{(\roman*)}]
\item The unique fixed point in $(0,2)$ is $\gamma = 0$.
\item $\gamma = 1$ is the curvature-optimal initial condition.
\item The orbit from $\gamma_0 = 1$ converges to $0$
  geometrically, with $\kappa(\alpha^n) < 1$ for $n \geq 1$.
\item For $\gamma_0 > 1$, the orbit crosses $\gamma = 1$
  (improving curvature) then continues past
  (degrading it): the curvature along the orbit is
  non-monotone.
\end{enumerate}
\end{proposition}



%======================================================================
\section{Analytical Consequences}\label{sec:consequences}
%======================================================================

The iterated curvature bounds have concrete analytical
content: they control the decay of the stable semigroup
at each energy level and yield functional inequalities
that interpolate between the process-specific behavior
at low chaos depth and universal Gaussian behavior at
high chaos depth.


\subsection{The energy hierarchy}

\begin{definition}[$k$-th energy functional]
\label{def:energy-k}
For $f \in \mathcal{T}_N^+$ and $k \geq 1$:
\begin{equation}\label{eq:Ek}
  \mathcal{E}_k(f)
  := \int_\T \Gamma_k(f,f)\,d\mu
  = \sum_{n=1}^N |a_n|^2\,n^{2Hk},
\end{equation}
where the second equality follows from orthogonality of
$e^{i(m-n)x}$ under $\mu = dx/(2\pi)$.
\end{definition}

\begin{remark}[Sobolev interpretation]
\label{rem:sobolev}
The $k$-th energy is the squared Sobolev-type norm at
fractional order $Hk = \gamma k/2$:
$\mathcal{E}_k(f) = \|f\|_{\dot{H}^{Hk}(\T)}^2$.
The hierarchy $\mathcal{E}_1, \mathcal{E}_2, \ldots$
measures progressively higher regularity.
\end{remark}

\begin{proposition}[Iterated spectral inequality]
\label{prop:iterated}
For all $k \geq 1$:
\begin{equation}\label{eq:iterated}
  \kappa^{(k)}(\gamma,N)\,\mathcal{E}_{k-1}(f)
  \;\leq\; \mathcal{E}_k(f)
  \;\leq\; K^{(k)}(\gamma,N)\,\mathcal{E}_{k-1}(f),
\end{equation}
where
$\kappa^{(k)} = \lambda_{\min}(R_H^{\circ k}, R_H^{\circ(k-1)})$
and
$K^{(k)} = \lambda_{\max}(R_H^{\circ k}, R_H^{\circ(k-1)})$.
Iterating the lower bound:
\begin{equation}\label{eq:chain}
  \mathcal{E}_k(f)
  \;\geq\; \Bigl(\prod_{j=2}^k \kappa^{(j)}\Bigr)
  \cdot \mathcal{E}_1(f).
\end{equation}
At $\gamma = 1$:
$\prod_{j=2}^k \kappa_1^{(j)} = 1$ (the chain is sharp
at the first mode) and
$\prod_{j=2}^k K^{(j)} = \prod_{j=2}^k (2N-1)(3N^2-3N+1)
\cdots \to \infty$ (the maximum grows polynomially in $N$
at each step).
\end{proposition}

\begin{proof}
Inequality~\eqref{eq:iterated} follows from
$\Gamma_k \geq \kappa^{(k)}\Gamma_{k-1}$ and
$\Gamma_k \leq K^{(k)}\Gamma_{k-1}$ after integration.
The chain~\eqref{eq:chain} is the telescope.
At $\gamma = 1$: $\kappa_1^{(j)} = 1$ for all $j$
(Corollary~\ref{cor:stability}).
\end{proof}


\subsection{Semigroup energy decay}

The stable semigroup $P_t = e^{tL_\gamma}$ acts as
$P_t(e^{inx}) = e^{-n^\gamma t}\,e^{inx}$.

\begin{lemma}[Energy dissipation identity]
\label{lem:dissipation}
For $k \geq 1$ and $f \in \mathcal{T}_N^+$:
\begin{equation}\label{eq:dissipation}
  \frac{d}{dt}\,\mathcal{E}_k(P_t f)
  = -2\,\mathcal{E}_{k+1}(P_t f).
\end{equation}
Each energy level is dissipated by the next.
\end{lemma}

\begin{proof}
From the $\Gamma$-recursion~\eqref{eq:Gammak-def}:
$\Gamma_k(f, L_\gamma f) =
\frac{1}{2}L_\gamma\Gamma_k(f,f) - \Gamma_{k+1}(f,f)$.
Integrating over $\T$ and using
$\int L_\gamma g\,d\mu = 0$:
\[
  \int_\T \Gamma_k(f, L_\gamma f)\,d\mu
  = -\mathcal{E}_{k+1}(f).
\]
Since $\frac{d}{dt}P_t f = L_\gamma P_t f$:
\[
  \frac{d}{dt}\mathcal{E}_k(P_t f)
  = 2\int_\T \Gamma_k(P_t f, L_\gamma P_t f)\,d\mu
  = -2\,\mathcal{E}_{k+1}(P_t f).
  \qedhere
\]
\end{proof}

\begin{theorem}[Semigroup energy decay hierarchy]
\label{thm:decay}
For $k \geq 1$ and $f \in \mathcal{T}_N^+$:
\begin{equation}\label{eq:decay}
  \mathcal{E}_k(P_t f)
  \;\leq\; e^{-2\kappa^{(k+1)}(\gamma,N)\,t}\,
  \mathcal{E}_k(f).
\end{equation}
The $k$-th energy decays exponentially at rate
$2\kappa^{(k+1)}$.
\end{theorem}

\begin{proof}
By Lemma~\ref{lem:dissipation} and the curvature bound
$\mathcal{E}_{k+1} \geq \kappa^{(k+1)}\mathcal{E}_k$:
\[
  \frac{d}{dt}\mathcal{E}_k(P_t f)
  = -2\,\mathcal{E}_{k+1}(P_t f)
  \leq -2\kappa^{(k+1)}\,\mathcal{E}_k(P_t f).
\]
Gr\"onwall's inequality gives~\eqref{eq:decay}.
\end{proof}

\begin{remark}[The standard case $k = 1$]
\label{rem:standard}
At $k = 1$: the Dirichlet energy $\mathcal{E}_1(P_t f)$
decays at rate $2\kappa^{(2)} = 2\kappa(\gamma,N)$, where
$\kappa(\gamma,N)$ is the Bakry--\'Emery constant from
Part~I.  At $\gamma = 1$: $\kappa = 1$, giving
$\mathcal{E}_1(P_t f) \leq e^{-2t}\,\mathcal{E}_1(f)$.
This is optimal: the mode $f = e^{ix}$ has
$\mathcal{E}_1(e^{ix}) = 1$ and
$\mathcal{E}_1(P_t e^{ix}) = e^{-2t}$.
\end{remark}


\subsection{Asymptotic Gaussian universality}

\begin{corollary}[Universal high-chaos decay rate]
\label{cor:universality}
For all $\gamma \in (0,2)$ and $\varepsilon > 0$,
there exists $k_0 = k_0(\gamma, N, \varepsilon)$ such that
for all $k \geq k_0$ and all $f \in \mathcal{T}_N^+$:
\begin{equation}\label{eq:universality}
  \mathcal{E}_k(P_t f)
  \;\leq\; e^{-2(1-\varepsilon)\,t}\,\mathcal{E}_k(f).
\end{equation}
That is, the decay rate at high energy levels
approaches $e^{-2t}$---the Gaussian rate---for all
stable generators, regardless of $\gamma$.
\end{corollary}

\begin{proof}
By Theorem~\ref{thm:regularization}:
$\kappa^{(k+1)} \geq 1 - C_N(\rho^*)^k \geq 1-\varepsilon$
for $k$ sufficiently large.
Apply Theorem~\ref{thm:decay}.
\end{proof}

This is the analytical content of chaos regularization:
every $\gamma$-stable semigroup, regardless of its
low-chaos curvature behavior, eventually decays at the
universal Gaussian rate at high energy levels.
The index $k_0$ depends on $\gamma$ through $\rho^*$:
near $\gamma = 1$ ($H \approx 1/2$), $\rho^*$ is
small and the Gaussian rate is reached quickly;
far from $\gamma = 1$, more chaos levels are needed.


\subsection{Exact rates at $\gamma = 1$}

At the Cauchy point, the eigenvalue formula
(Theorem~\ref{thm:eigenvalues}) gives exact decay
rates for each Fourier mode.

\begin{proposition}[Mode-resolved decay]
\label{prop:mode-decay}
At $\gamma = 1$, for each mode $f = e^{inx}$
($n = 1,\ldots,N$):
\begin{equation}\label{eq:mode-decay}
  \frac{\mathcal{E}_{k+1}(P_t e^{inx})}
       {\mathcal{E}_k(P_t e^{inx})}
  = n \qquad \text{for all } t \geq 0.
\end{equation}
More generally, for any $f \in \mathcal{T}_N^+$:
\begin{equation}\label{eq:sandwich}
  1\;\leq\;
  \frac{\mathcal{E}_{k+1}(P_t f)}{\mathcal{E}_k(P_t f)}
  \;\leq\; N,
\end{equation}
and both bounds are attained (by $e^{ix}$ and $e^{iNx}$
respectively).
\end{proposition}

\begin{proof}
For $f = e^{inx}$:
$\mathcal{E}_k(P_t e^{inx}) = n^k\,e^{-2nt}$.
The ratio is $n^{k+1}/n^k = n$, independent of $t$.
The sandwich follows from
$\min_n n = 1$ and $\max_n n = N$.
\end{proof}

\begin{remark}[The complete picture at $\gamma = 1$]
\label{rem:complete}
Combining Proposition~\ref{prop:mode-decay} with the
eigenvalue formula: the semigroup acts on the $k$-th
energy as multiplication by $e^{-2n^\gamma t}$
mode-by-mode, and the curvature
$\kappa^{(k+1)} = \min_j \kappa_j^{(k+1)}$ controls the
worst-case (lowest-mode) decay.  At $\gamma = 1$:
$\kappa_1^{(k+1)} = 1$ for all $k$, giving the uniform
rate $e^{-2t}$.  Higher modes decay faster:
$e^{inx}$ has decay rate $e^{-2nt}$, with the $k$-th
energy ratio being $n$ regardless of $k$.
The Bakry--\'Emery hierarchy at $\gamma = 1$ is thus
completely transparent: curvature is $1$ at all
chaos levels, and each mode decays at its own
spectral rate.
\end{remark}


%======================================================================
\section{Discussion and Open Problems}\label{sec:discussion}
%======================================================================

\subsection{The correspondence dictionary}

The results of Parts~I and~II extend the stable--fBM
dictionary:
\[
\begin{array}{c|c}
\textbf{Stable process / $\Gamma$-calculus}
& \textbf{fBM / operator derivative} \\[4pt]
\hline
\rule{0pt}{14pt}
\Gamma(f,f) \text{ (carr\'e du champ)}
  & \|D_{B^H}\Phi_f\|^2_{\HH}
  \text{ (energy norm)} \\
\Gamma_k(f,f) \text{ (iterated CdC)}
  & \|\iota_k(\mathbf{b})\|^2_{\HH^{\otimes k}}
  \text{ (tensor norm)} \\
R_H^{\circ k} \text{ (Hadamard power)}
  & \text{Gram matrix of diagonal } k\text{-tensors} \\
\text{$k$-th curvature } \kappa^{(k)}
  & \text{contraction ratio } \iota_k / \iota_{k-1} \\
\text{Eigenvalues at } \gamma = 1
  & \|\delta_j^{(k)}\|^2 / \|\delta_j^{(k-1)}\|^2 \\
\text{Cholesky } M_k = LD_kL^T
  & \text{increment-tensor expansion} \\
\text{Stability } \kappa = 1
  & \|e_1^{\otimes k}\| = 1 \text{ (unit first increment)} \\
\text{Curvature deficit}
  & \text{increment correlation at } H \neq 1/2 \\
\text{Subordination } \gamma \mapsto \alpha\gamma
  & H \mapsto \alpha H \text{ (energy space change)} \\
\text{Chaos regularization}
  & \text{off-diagonal decay at rate } (\rho^*)^{k-1}
\end{array}
\]


\subsection{Relation to Part~I}

Theorem~F strengthens Part~I's phase transition
by providing the strict inequality
$\kappa(\gamma,N) < 1$ for all $\gamma \neq 1$ and
$N \geq 2$ via perturbation theory.
The full curvature landscape, combining both papers, is:
\begin{enumerate}[label=(\roman*)]
\item For $\gamma \in (0,1]$:
  $\kappa(\gamma) = \lim_N \kappa(\gamma,N)$ exists and
  satisfies $\frac{1}{2} \leq \kappa(\gamma) \leq 1$.
  (For $\gamma < 1$, the lower bound $\frac{1}{2}$ is
  conditional on the Z-matrix property of
  $R_H^{-1}R_H^{\circ 2}$; see Part~I, Lemma~5.5.)
\item At $\gamma = 1$: $\kappa(1) = 1$ (unique maximum).
\item For $\gamma \in (1,2)$:
  numerical evidence (Part~I, Conjecture~5.8)
  suggests $\kappa(\gamma) > 0$, but
  the Z-matrix proof strategy fails for $H > 1/2$.
\end{enumerate}


\subsection{Open problems}

\begin{enumerate}[label=(\roman*)]
\item \textbf{Analytical proof of global optimality.}
  Theorem~\ref{thm:optimality} is verified numerically.
  An analytical proof of the strict negativity of the
  second-order perturbation coefficient $c_N > 0$
  would make the result fully rigorous.
  The correlation factorization of
  Section~\ref{sec:regularization} suggests an approach:
  prove that the perturbation $\rho^{\circ 2}(H)$ at
  $H = 1/2 + \delta$ loses positive definiteness
  relative to $\rho(H)$ at second order in $\delta$,
  using the explicit formula~\eqref{eq:rho}.

\item \textbf{Sharp rate of chaos regularization.}
  Theorem~\ref{thm:regularization} gives the rate
  $O((\rho^*)^{k-1})$ with a Gershgorin-based constant
  $C_N = 2(N{-}1)\,N^{2H}$.  Is the rate $\rho^*$ sharp,
  or does the true rate involve a smaller quantity
  (e.g., the spectral radius of $\rho - I$ restricted
  to the eigenspace of $\Delta^2$ near eigenvalue $1$)?
  The constant $C_N$ grows with $N$; understanding its
  sharpness is related to the $N \to \infty$ limit.

\item \textbf{Iterated operator derivatives.}
  The tensor realization (Theorem~\ref{thm:tensor}) uses
  the diagonal lift $\iota_k$ to connect $\Gamma_k$ to
  $\HH^{\otimes k}$.  This should correspond to the
  $k$-th iterated operator derivative
  $D_{B^H}^k$, which maps the $k$-th Wiener chaos
  isometrically (up to $k!$) onto $\HH^{\odot k}$
  (symmetric tensors).  Making this correspondence
  precise---showing that the diagonal lift is the
  restriction of the iterated operator derivative
  to the mode vectors---is an open problem.

\item \textbf{Drift at higher chaos levels.}
  Part~I showed that the cosine drift acts as a scalar
  shift at level $k = 2$ for $\gamma = 1$.
  Does this scalar property persist at higher levels?
  In the tensor framework, the drift correction
  at level $k$ should be
  $\frac{\omega^2}{2}\cos(x) \cdot M_{k-1}$,
  giving eigenvalues
  $\kappa_j^{(k)}(x) = \kappa_j^{(k)}
  + \frac{\omega^2}{2}\cos x$.

\item \textbf{Multi-dimensional torus.}
  On $\T^d$, the Hadamard power hierarchy produces
  anisotropic curvature tensors.  The Cholesky
  structure generalizes to tensor products of
  min-matrices, and the energy space becomes
  $\HH^{\otimes d}$.

\item \textbf{Higher-order curvature at $\gamma \neq 1$.}
  For $\gamma \neq 1$, are the eigenvalues of
  $(R_H^{\circ(k-1)})^{-1}R_H^{\circ k}$ computable
  in closed form?  The Cholesky structure
  $M_k = LD_kL^T$ is specific to $H = 1/2$;
  for general $H$, the increment tensors lose
  orthogonality and the factorization involves
  non-trivial Cholesky factors.
\end{enumerate}


%======================================================================
\section*{Acknowledgments}
%======================================================================
AI tools were used as interactive assistants for mathematical
exploration, numerical verification, and drafting. The author
is responsible for all mathematical content and exposition.


%======================================================================
\begin{thebibliography}{20}

\bibitem{BakryGentilLedoux2014}
D.~Bakry, I.~Gentil, and M.~Ledoux.
\newblock \emph{Analysis and Geometry of Markov Diffusion Operators}.
\newblock Grundlehren Math. Wiss., vol.~348, Springer, 2014.

\bibitem{BoettcherSchillingWang2013}
B.~B\"ottcher, R.~Schilling, and J.~Wang.
\newblock \emph{L\'evy Matters III: L\'evy-Type Processes}.
\newblock Lecture Notes in Math., vol.~2099, Springer, 2013.

\bibitem{Fontes2026a}
R.~Fontes.
\newblock Stochastic calculus as operator factorization.
\newblock \emph{arXiv preprint} arXiv:2601.09976 [math.PR], 2026.

\bibitem{Fontes2026b}
R.~Fontes.
\newblock Bakry--\'Emery curvature of the fractional Laplacian
  via fractional Brownian covariance.
\newblock \emph{Preprint}, 2026.

\bibitem{Fontes2026opderiv}
R.~Fontes.
\newblock The operator derivative in continuous stochastic calculus:
  a Hilbert energy space framework.
\newblock \emph{Preprint}, 2026.

\bibitem{Garofalo2019}
N.~Garofalo.
\newblock Fractional thoughts.
\newblock In \emph{New Developments in the Analysis of Nonlocal Operators},
  Contemp. Math., vol.~723, AMS, 2019.

\bibitem{Mishura2008}
Y.~Mishura.
\newblock \emph{Stochastic Calculus for Fractional Brownian Motion
  and Related Processes}.
\newblock Lecture Notes in Math., vol.~1929, Springer, 2008.

\bibitem{Schoenberg1938}
I.~J.~Schoenberg.
\newblock Metric spaces and positive definite functions.
\newblock \emph{Trans. Amer. Math. Soc.}, 44(3):522--536, 1938.

\bibitem{SpenerWeberZacher2020}
A.~Spener, F.~Weber, and R.~Zacher.
\newblock The fractional {L}aplacian has infinite dimension.
\newblock \emph{Comm. Partial Differential Equations}, 45(1):57--75, 2020.

\end{thebibliography}

\end{document}
